\documentclass[11pt]{article}
\usepackage[margin=1in]{geometry}
\usepackage{amsmath,amssymb}
\usepackage[hidelinks]{hyperref}

\title{Disproof of Conjecture 00000009614\\
\large The same Perron eigenvalue does not determine the dimension group}
\author{}
\date{}

\begin{document}
\maketitle

Write \(J_n\) for the \(n\times n\) all-ones matrix. Take the two
symmetric zero-one adjacency matrices
\[
 A=J_2,\qquad B=J_3-I_3.
\]
They define mixing shifts of finite type because \(A\) is positive
and \(B^2=J_3+I_3\) is positive. Each has row sum two, so the
all-ones vector is an eigenvector of eigenvalue two. Each eigenvalue
has absolute value at most the maximum row sum, giving
\(\lambda_A=\lambda_B=2\). Thus the generated real number fields
are both \(\mathbb Q\), and the two eigenvalues have the same unit
coset in that field, with ratio the unit \(1\).

\section*{The underlying dimension groups}

The underlying group \(G_T\) of the Krieger dimension group of an
adjacency matrix \(T\) is the stationary direct limit of
\(\mathbb Z^n\) under \(T^{\mathsf T}\)
\cite[Section~3]{matsumoto}. Both matrices here are symmetric,
so \(T^{\mathsf T}=T\). A vector at stage \(k\) represents a
class \([v,k]\), and two representatives agree exactly when their
images at some common later stage agree.

For \(A\), the map
\[
  [(a,b),k]\longmapsto (a+b)/2^k
\]
is an isomorphism onto \(\mathbb Z[1/2]\). The numerator sum
doubles under each transition, so the map is well-defined.
It is onto, and a vector of sum zero is killed by the next
transition, proving injectivity. For \(B\), the map
\([v,k]\longmapsto B^{-k}v\) embeds \(G_B\) into \(\mathbb Q^3\),
since \(B^{-1}=\frac12J_3-I_3\) over \(\mathbb Q\). Its image
contains the stage-zero copy of \(\mathbb Z^3\). Thus \(G_A\)
has rational rank one and \(G_B\) has rational rank three.
They cannot be isomorphic as abelian groups or as ordered groups.

\section*{Finite quotient verification}

The distinction is also witnessed by the quotient modulo three.
Reduction of each stage modulo three gives a surjection
\[
  G_T\longrightarrow
  \varinjlim\bigl(\mathbb F_3^n\xrightarrow{\bar T}
                 \mathbb F_3^n\xrightarrow{\bar T}\cdots\bigr).
\]
It is onto because each vector of residues has an integer lift.
A representative \([v,k]\) maps to zero precisely when some
\(\bar T^r\bar v=0\). Then \(T^r v=3w\) for an integer vector \(w\)
and \([v,k]=3[w,k+r]\). Conversely every element of \(3G_T\)
maps to zero. This identifies the finite-field limit with
\(G_T/3G_T\).

For \(\bar A\), one step moves each pair to the diagonal, and
on the diagonal the map doubles the coordinate. Doubling is
invertible in \(\mathbb F_3\), so the limit has three elements.
For \(\bar B\), the formula \(2J_3-I_3\) gives an inverse in
\(\mathbb F_3\). Every stage therefore identifies with the full
space \(\mathbb F_3^3\), and the limit has \(27\) elements.
An isomorphism of abelian groups induces an isomorphism on the
quotients modulo three, giving the same contradiction.

\section*{Lean verification}

The Lean~4.19.0 project forms the stationary limits as quotients
by eventual equality of stage-vector pairs. It proves the
representative correspondence for both finite-field systems
and verifies the complete enumerations of \(3\) and \(27\)
representatives. Its final theorem excludes a surjection from
the first limit to the second. The identification with the
integer dimension groups modulo three is proved above. Run
\verb|lake build| in \verb|lean/| to verify the formal results.

\begin{thebibliography}{1}
\bibitem{matsumoto}
K.~Matsumoto, \emph{Flow equivalence of topological Markov shifts and
Ruelle algebras}, New York J. Math. \textbf{27} (2021), 1375--1414.
\href{https://nyjm.albany.edu/j/2021/27-53v.pdf}{Article}.
\end{thebibliography}

\end{document}
